
\begin{proof}
Let $V$ be an $n$-element set and partition $V$ into disjoint sets $V = V_0 \sqcup V_1$ where $|V_0| = \floor{n/2}$, $|V_1| = \ceil{n/2}$.  Choose a distinguished vertex $v_0 \in V_0$, and let $V_0' = V_0 \setminus \{v_0\}$. 
Let $N = |V_0' \times V_1|$ and let $g : V_0' \times V_1 \rightarrow [N]$ be a bijection.  For every $x \in \{0,1\}^N$ we define a weighted graph $G_x = (V,w_x)$ where
\begin{itemize}
\item
$w_x(\{u,v\}) = \tau$ if $u,v \in V_0$ or $u,v \in V_1$,
\item
$w_x(\{u,v\}) = x(g(\{u,v\}))$ if $(u \in V_0', v \in V_1)$ or $(u \in V_1, v \in V_0')$,
\item
$w_x(\{u,v\}) = 0$ otherwise.
\end{itemize}
Let $k = \tau (\floor{n/2}- 1)$.  In the following $\Delta_{G_x}(\cdot)$ will always be with respect to $G_x$ and we drop the subscript.
For any $x$ it holds that $w_x(\Delta(V_0)) = |x|$ and $w_x(\Delta(\{v_0\})) = k$.
Now consider any $x$ and a subset $\emptyset \neq Y \subsetneq V$ different from $V_0$ or $V_1$.  We can prove that $w_x(\Delta(Y)) \ge k$.  To this end note that either 
$\emptyset \neq Y \cap V_0 \subsetneq V_0$ or $\emptyset \neq Y \cap V_1 \subsetneq V_1$.  First assume that the former is the case.  Then
\[
w_x(\Delta(Y))
= \sum_{u \in Y, v \notin Y} w_x(\{u,v\})
\ge \sum_{u \in Y \cap V_0, v \in V_0 \backslash Y} w_x(\{u,v\})
\ge k,
\]
as $k$ is the weight of a minimum cut in the complete weighted graph over $\floor{n/2}$ nodes with all edge weights $\tau$.  If instead $\emptyset \neq Y \cap V_1 \subsetneq V_1$ 
then a similar argument shows that $w_x(\Delta(Y)) \ge \tau (\ceil{n/2}-1) \ge k$.

Thus if $|x| < k$ then $\Delta(V_0)$ will be the unique minimum cut of $G_x$, and the weight of a minimum cut in $G_x$ will be $w_x(\Delta(V_0)) = |x|$. On the other hand, if 
$|x| > k$ then the weight of a minimum cut in $G_x$ will be $k$, which is realized by the star cut $\Delta(\{v_0\})$ (and potentially other cuts in $G_x$) but 
not by $\Delta(V_0)$ as $w_x(\Delta(V_0)) = |x| > k$.

Let $\Scal = \{x \in \{0,1\}^N : |x| \in \{k-1,k+1\} \}$ and $\Gcal = \{G_x : x \in \Scal\}$.  
Suppose there was a $T$ query algorithm in the adjacency matrix model that for any $G_x \in \Gcal$ with probability at least $2/3$ output the weight of a minimum cut in $G_x$.
If the output is $<k$ then we know that $|x| = k-1$ and if the output is $k$ then we know that $|x| = k+1$.  Moreover, any query to the adjacency matrix of $G_x$ can be simulated 
by a query to $x$, thus such an algorithm gives a $T$ query algorithm to determine if $|x| = k-1$ or $|x| = k+1$ when we are promised one of these is the case.  Since $\tau \le (\floor{n/2}-1)/2$ 
we have $2k+1 \le N$ and therefore we may apply Theorem \ref{thm:thresh} to obtain $T \in \Omega(\sqrt{N k}) = \Omega(n^{3/2}\sqrt{\tau})$.

Similarly, a $T$ query algorithm in the adjacency matrix model that for any $G_x \in \Gcal$ with probability at least $2/3$ outputs the shores of a cut realizing the minimum weight 
also implies a $T$ query algorithm to determine if $|x| = k-1$ or $|x| = k+1$.  In this case, if $|x| = k-1$ then the output must be $(V_0, \overline{V}_0)$ as these are the shores of the unique 
minimum cut in $G_x$.  On the other hand, if $|x| = k+1$ then $(V_0, \overline{V}_0)$ is not a correct output.  Thus the output of the algorithm lets us determine with probability at least $2/3$ if 
$|x| = k-1$ or $|x| = k+1$ and we again have $T \in \Omega(\sqrt{N k}) = \Omega(n^{3/2}\sqrt{\tau})$.
\end{proof}

\subsection{Adjacency array model}
Given adjacency array access to a graph with edge-weight ratio $\tau$, we showed an upper bound of $\tO(\sqrt{m n \tau})$ on the quantum query complexity of computing the weight of a minimum cut.
In this section we prove two distinct lower bounds, each of which is tight in a specific regime.
First we show that for any $\tau > 1$ there exists a family of dense graphs on $n$ vertices with edge-weight ratio $\tau$ for which computing the weight of a minimum cut requires $\Omega(n^{3/2})$ queries to the adjacency array.
This shows that the adjacency array upper bound of Theorem~21 of the source is tight for dense weighted graphs with constant (but non-unit) edge-weight ratio.
Secondly and using a different approach, for any $1 \leq \tau \in O(n)$ we prove an $\Omega(\tau n)$ lower bound for a family of dense graphs with edge-weight ratio $\tau$.
This shows that we cannot get a quantum speedup when $\tau \in \Omega(n)$.

\subsubsection{Constant edge-weight ratio}

For the first bound we first need a claim about the minimum cuts of a complete weighted bipartite graph.
\begin{claim}
\label{clm:complete_bipartite}
Let $n \ge 8$ be a multiple of $4$ and $G=(L \sqcup R,w)$ be a weighted bipartite graph with bipartition $L,R$ where $|L| = 3n/4, |R| = n/4$. Further suppose that 
for every $x \in L, y \in R$ it holds that $w(\{x,y\}) \ge 1$.   Then any cut of $G$ that is not of the form $\Delta_G(\{x\})$ for $x \in L$ has weight at least $n/2$.  
\end{claim}

\begin{proof}
First consider a star cut $\Delta_G(\{y\})$ for $y \in R$.  This has weight at least $3n/4$, since this is the degree of $y$ and all edges have weight at least 1.  

It now remains to show the claim holds for non-star cuts.  Consider a general non-star cut with shore $X \cup Y$ with $X \subseteq L, Y \subseteq R$.  Let 
$k = |X|, \ell = |Y|$.  As it is a non-star cut we have $k + \ell \ge 2$.  By complementing as needed we may also assume that $k \le 3n/8$.  We also have 
the obvious constraints that $\ell \le n/4$ and $k,\ell \ge 0$.  

As $G$ is a complete weighted bipartite graph with every edge weight at least one we have
\[
w(\Delta_G(X \cup Y)) \ge k(n/4-\ell) + \ell(3n/4-k) \enspace .
\]
As $k \le 3n/8$ the term $\ell (3n/4-k)$ is greater than $n/2$ whenever $\ell \ge 2$.  Thus we can focus on $\ell \in \{0,1\}$.  
If $\ell=0$ then $k\ge 2$ and so the weight of the cut is at least $k(n/4)=n/2$ as desired.  If $\ell = 1$ then the weight of the 
cut is $k(n/4-1) + 3n/4-k$ which is always at least $3n/4$ as long as $n \ge 8$.  
\end{proof}

This claim means that if $\min_{x \in L} w(\Delta_G(\{x\})) < n/2$ then this value will be the weight of a minimum cut in $G$.  
We can leverage this to show a lower bound as follows.
In the next proof, for a function $f: \{0,1\}^n \rightarrow \{0,1\}$ we will use $Q_{1/3}(f)$ to denote the quantum query complexity of computing $f$ with error at most $1/3$.
